% ============================================================
%  Chapter 8 — Conclusion
% ============================================================
\section{Conclusion}
\label{ch:conclusion}

\subsection{Summary of Findings}
\label{sec:conc:summary}

This thesis investigated the operational planning problem of a Baltic Sea
\RoRo{} vessel subject to the IMO Carbon Intensity Indicator (\CII{})
regulation and EU Emissions Trading System (\ETS{}) carbon pricing.  A
mixed-integer linear programme (\MILP{}) was formulated to minimise
voyage cost by optimising discrete speed choices across a multi-arc
service cycle, subject to per-arc time budgets, a fleet-level \CII{}
compliance constraint, and the RoRo-specific \CII{} formula based on
gross tonnage.  Six scenarios were evaluated against the EU MRV dataset
for M/V Obbola, a purpose-built forest-products \RoRo{} vessel operated
by \SCA{} Logistics in the Baltic Sea.

Four main findings emerge.

\paragraph{Finding 1: The published schedule, not the CII target, is the
binding operational constraint under the 2026 regulatory framework.}
In Scenario S3, the cost-optimal solution achieves a Rating~A \CII{}
of 10.2705~\gGTnm{} --- 33.9\,\% below the 2026 Rating~C threshold ---
not because of active \CII{} management but because the time budgets on
four bottleneck arcs force the vessel to sail at speeds (14--15~kn) that
already satisfy the regulation with substantial headroom.  Speed reduction
below these minima is precluded by the schedule.

\paragraph{Finding 2: A 2030 Rating~B CII target is structurally
infeasible under the current timetable.}
The analytically derived minimum achievable \CII{} of 10.3187~\gGTnm{}
exceeds the 2030 Rating~B threshold of 9.9291~\gGTnm{} by 3.9\,\%.
No combination of speed choices satisfying the published time budgets
can achieve Rating~B; the infeasibility is structural, arising from
four specific bottleneck arcs.  This constitutes an original contribution
of the thesis.

\paragraph{Finding 3: Schedule flexibility yields significant emission
and cost reductions.}
Replacing per-arc time budgets with a global cycle-time cap (Scenario S5)
reduces the attained \CII{} from 10.2705 to 8.7116~\gGTnm{} (a 15.2\,\%
improvement) and generates meaningful fuel cost savings relative to the
fixed-schedule optimum.  This quantifies the \emph{value of schedule
flexibility} and suggests that operator negotiations over port windows are
a more tractable emission-reduction lever than on-board technology for
this vessel type.

\paragraph{Finding 4: The RoRo CII formula assigns responsibility
differently from transport-work-based metrics.}
By using gross tonnage in the denominator rather than transport work
(capacity $\times$ distance), the \CII{} formula for \RoRo{} vessels
is insensitive to cargo utilisation.  This design choice rewards
operators for reducing fuel consumption per GT-mile but does not provide
incentives for maximising cargo fill.  The finding has implications
for how the regulation should be interpreted and potentially reformed
for the \RoRo{} sector.

\subsection{Contributions}
\label{sec:conc:contributions}

This thesis makes the following scholarly and practical contributions:

\begin{enumerate}
  \item A \MILP{} formulation for cost-optimal speed assignment over a
    multi-arc \RoRo{} service cycle under a linearised \CII{} compliance
    constraint, extended with a flexible-schedule variant.

  \item An analytical minimum-CII function that derives the lowest
    structurally achievable \CII{} under a fixed timetable, enabling
    direct detection of regulatory infeasibility without enumeration.

  \item A documented case study of M/V Obbola on the SCA Baltic Sea
    route, including model validation against EU MRV reported emissions
    (1.2\,\% \CII{} error), demonstrating the applicability of the
    framework to a real short-sea \RoRo{} operation.

  \item The identification of structural \CII{} infeasibility for the
    2030 Rating~B target on the current SCA timetable, alongside a
    quantification of the required schedule relaxation or technological
    improvement needed to close the gap.

  \item A quantification of the value of schedule flexibility, expressed
    as a reduction in attained \CII{} and voyage fuel cost, providing
    a tangible metric for schedule-renegotiation decisions.
\end{enumerate}

\subsection{Limitations and Future Work}
\label{sec:conc:future}

Several limitations of the present work point to avenues for future
research.

The most immediately actionable extension is to incorporate stochastic
cargo weights and fuel prices into the optimisation, replacing the
deterministic single-value parameters with distributions fitted to
historical observations.  This would yield robust speed policies that
hedge against the dominant sources of voyage cost variability.

A second extension concerns the post-2026 reduction factor.  When the
\IMO{} MEPC adopts confirmed annual reduction factors beyond 2026
(expected at MEPC 83/84), the scenario set should be re-run to assess
whether Rating~C targets also become infeasible before 2030, and at
what reduction level the infeasibility threshold is crossed.

Third, the thesis models a single vessel on a single route.  A
multi-vessel, multi-route extension --- covering the full SCA fleet
operating in the Baltic and North Sea --- would allow fleet-level
\CII{} averaging strategies to be evaluated, which may provide
compliance flexibility not available at the individual vessel level.

Fourth, the model does not include alternative fuel options.  As the
regulatory pressure beyond 2030 intensifies and as LNG, methanol, and
eventually ammonia become commercially available for Baltic Sea
operators, a comparative techno-economic model incorporating fuel
switching could identify cost-effective pathways to decarbonisation.

Finally, the non-monotonicity in the SCA fuel polynomial at 12--13~kn
deserves further investigation.  If confirmable from engine test data,
it could be exploited operationally; if it is a model artefact, the
polynomial should be recalibrated with additional sea-trial observations.

\subsection{Closing Remarks}
\label{sec:conc:closing}

The \CII{} regulation represents the maritime sector's most concrete
near-term policy instrument for curbing operational emissions.  This
thesis shows that its interaction with fixed-schedule short-sea services
is more complex than the regulation's design may suggest: in the near
term, operators are already compliant largely by accident of their
timetables; in the medium term, the pathway to stricter targets requires
not just slower sailing but a renegotiation of the fundamental service
architecture.

The contribution of rigorous optimisation modelling to this debate is
not to recommend a single speed profile but to make the constraints
explicit, to identify where the regulatory and operational systems are
misaligned, and to quantify the value of actions --- such as schedule
flexibility --- that might otherwise be overlooked.  It is hoped that
the framework developed here provides a useful foundation for further
work on cost-effective decarbonisation in the European short-sea shipping
sector.
